\chapter*{CAPÍTULO IV: ESPECIFICACIONES TÉCNICAS DE MONTAJE ELECTROMECÁNICO}
\addcontentsline{toc}{chapter}{CAPÍTULO IV: ESPECIFICACIONES TÉCNICAS DE MONTAJE ELECTROMECÁNICO}

\section{NORMAS Y PROCEDIMIENTOS GENERALES}

El montaje de las obras electromecánicas se ejecutará en concordancia con la Norma DGE \textit{Especificaciones Técnicas para el Montaje de Líneas Aéreas y Redes de Distribución} (R.D. N.° 026-2003-EM/DGE), las prescripciones del Código Nacional de Electricidad - Suministro 2011, las normas de seguridad ocupacional en electricidad (R.M. N.° 111-2013-MEM/DM - RESESATE) y los planos de detalle aprobados.

El Contratista ejecutor dispondrá en obra de un Ingeniero Mecánico Electricista o Electricista colegiado y habilitado como Residente de Obra, así como de personal técnico calificado con experiencia acreditada en líneas de media tensión y redes de distribución urbana.

\section{TRAZO, ESTACADO Y REPLANTEO TOPOGRÁFICO}

\begin{enumerate}
  \item \textbf{Georreferenciación y Catastro:} Se utilizará como control primario la red geodésica del Catastro Urbano Consolidado de Limbani (WGS 84, UTM Zona 19 Sur). La ubicación de cada poste y subestación deberá respetar estrictamente las líneas municipales, anchos de vía y servidumbres de paso existentes en el catastro.
  \item \textbf{Estacado:} Se colocarán estacas de madera cepillada con pintura fluorescente en el punto exacto del eje de cada poste, numeradas correlativamente según la codificación de planos.
  \item \textbf{Verificación de vanos y perfiles:} Se comprobarán las cotas de terreno, desniveles críticos en la quebrada del Río Limbani y distancias mínimas de seguridad verticales y horizontales respecto a edificaciones y vegetación circundante.
\end{enumerate}

\section{EXCAVACIÓN DE HOYOS Y ZANJAS}

\subsection{Dimensiones de Excavación para Postes}
La profundidad de enterramiento de los postes de concreto armado centrifugado se rige por la fórmula reglamentaria:
$$H_e = \frac{H_t}{10} + 0.50\text{ m}$$

Donde $H_t$ es la longitud total del poste:
\begin{itemize}
  \item \textbf{Postes de $13.0\text{ m}$ (Media Tensión / SAB):} $H_e = \frac{13.0}{10} + 0.50 = \mathbf{1.80\text{ m}}$. Diámetro excavación: $0.80\text{ m}$.
  \item \textbf{Postes de $12.0\text{ m}$ (Media Tensión alineamiento):} $H_e = \frac{12.0}{10} + 0.50 = \mathbf{1.70\text{ m}}$. Diámetro excavación: $0.75\text{ m}$.
  \item \textbf{Postes de $9.0\text{ m}$ (Baja Tensión / Mixto):} $H_e = \frac{9.0}{10} + 0.50 = \mathbf{1.40\text{ m}}$. Diámetro excavación: $0.65\text{ m}$.
  \item \textbf{Postes de $8.0\text{ m}$ (Baja Tensión pasajes):} $H_e = \frac{8.0}{10} + 0.50 = \mathbf{1.30\text{ m}}$. Diámetro excavación: $0.60\text{ m}$.
\end{itemize}

\subsection{Excavación en Terreno Rocoso y Cemento Ciclópeo}
En zonas con afloramiento de roca compacta en las laderas de Limbani, las excavaciones se realizarán con martillo neumático o demoledor eléctrico, prohibiéndose el uso de dinamita cerca a viviendas habitadas. En todos los apoyos en ladera o esquinas sometidas a tracción, la base del poste se fundará con un solado de concreto simple de $0.15\text{ m}$ de espesor ($f'_c = 100\text{ kg/cm}^2$) y un macizo de anclaje de concreto ciclópeo ($f'_c = 140\text{ kg/cm}^2 + 30\%$ de piedra mediana).

\section{IZAJE, VERTICALIZACIÓN Y CIMENTACIÓN DE POSTES}

\begin{itemize}
  \item El izaje se ejecutará mediante pluma hidráulica sobre camión grúa en vías transitables, o mediante trípode mecánico con tecle de cadena (mástil de madera/acero) y vientos guía en los pasajes estrechos y empinados del casco urbano.
  \item Se verificará el aplomado con plomada topográfica en dos planos perpendiculares, asegurando una desviación máxima no mayor a $2.0\text{ mm}$ por metro de altura sobre el nivel de suelo.
  \item La orientación de los orificios de la cima debe ser perfectamente perpendicular a la dirección de la traza para los armados de alineamiento, o alineada con la bisectriz en ángulos de desvío.
  \item El relleno se compactará en capas sucesivas de $0.20\text{ m}$ de espesor con apisonador mecánico o pisón manual de mano dura, adicionando agua óptima de compactación.
\end{itemize}

\section{MONTAJE DE ARMADOS Y FERRETERÍA}

\begin{itemize}
  \item Todos los pernos, tuercas y arandelas deberán ajustarse con torquímetro al par especificado por el fabricante, aplicando arandelas curvas y de presión para evitar el aflojamiento por vibración eólica.
  \item En los postes de CAC se utilizarán abrazaderas partidas de dos semicilindros de plancha de acero galvanizado para evitar la concentración de esfuerzos puntuales sobre el fuste del poste.
  \item Los aisladores de espiga (pin) se fijarán rígidamente con espigas de acero forjado con cabeza de plomo roscado. Los aisladores poliméricos de suspensión deberán quedar alineados en el plano de tracción resultante.
\end{itemize}

\section{TENDIDO, TESADO Y FLECHADO DE CONDUCTORES EN MEDIA TENSIÓN}

\begin{enumerate}
  \item \textbf{Cuidado del Conductor AAAC:} Las bobinas de conductor se montarán sobre gatos mecánicos portabobinas con freno regulable. El conductor se tenderá utilizando poleas de aluminio recubiertas de neopreno colgadas en cada cruceta, evitando cualquier contacto del conductor con piedras, rocas o aristas filosas que provoquen entalladuras o rayas superficiales.
  \item \textbf{Control de Tracción y Flecha:} El tensado se realizará con tirfor o cabestrante mecánico y dinamómetro calibrado. Se medirá la temperatura del conductor mediante termómetro de contacto antes del amarre definitivo.
  \item \textbf{Regulación de Flecha:} Las flechas se comprobarán visualmente con crucetas de nivelación o estación total en el vano testigo, ajustando el valor exacto a las tablas de tensado para la temperatura ambiente de la jornada de trabajo.
  \item \textbf{Ataduras y Amarres:} En aisladores de pin se realizarán ataduras preformadas de aleación de aluminio sobre cinta de armar (armor tape) de aluminio para amortiguar esfuerzos cortantes.
\end{enumerate}

\section{TENDIDO Y ANCLAJE DE CABLES AUTOPORTANTES CAAI EN BAJA TENSIÓN}

\begin{itemize}
  \item El cable autoportante se desenrollará tirando del portador neutro mensajero (NA 25), manteniendo las fases protegidas de tracción mecánica directa.
  \item Se emplearán poleas de rodillos de nilón en cada estructura de paso.
  \item El anclaje terminal y de retención se realizará exclusivamente sobre el conductor neutro mensajero mediante grapas cónicas de aleación de aluminio de cuatro cuñas autoblocantes montadas sobre cáncamos o pernos de ojo pasantes.
  \item En apoyos de paso recto se suspenderá el neutro portador mediante grapas de suspensión con fusible mecánico de rotura controlada para protección ante caídas de árboles o postes.
  \item La flecha máxima del haz en servicio no deberá exceder el $1.5\%$ de la longitud del vano a la temperatura máxima de $35^\circ\text{C}$.
\end{itemize}

\section{MONTAJE DE SUBESTACIONES AÉREAS BIPOSTE (SAB)}

\subsection{Estructura Biposte y Plataforma}
\begin{itemize}
  \item La SAB se conformará con dos postes de CAC de $13.0\text{ m} / 400\text{ daN}$ separados una distancia entre ejes de $2.00\text{ m}$ a $2.20\text{ m}$.
  \item La plataforma de soporte del transformador estará constituida por perfiles estructurales UPN de acero galvanizado en caliente de $100 \times 50 \times 6\text{ mm}$, anclados a los postes a una altura de $4.50\text{ m}$ sobre el nivel del suelo mediante abrazaderas de cuatro pernos.
  \item Se dispondrá de baranda perimétrica de seguridad y piso de perfiles antideslizantes para las labores de inspección.
\end{itemize}

\subsection{Montaje del Transformador}
\begin{itemize}
  \item El transformador se izará mediante camión grúa provisto de eslingas de carga textil homologadas en los cuatro cáncamos de izaje del tanque.
  \item Se apoyará rígidamente sobre topes antivibratorios y se apernará a las vigas de la plataforma para impedir cualquier desplazamiento ante sismos.
\end{itemize}

\subsection{Protección y Maniobra en Media Tensión}
\begin{itemize}
  \item En la cruceta superior (altura mínima $8.50\text{ m}$) se instalarán tres (03) seccionadores cut-out unipolares de $15\text{ kV} / 100\text{ A}$ y tres (03) pararrayos de óxido metálico de $12\text{ kV} / 10\text{ kA}$.
  \item La conexión entre la línea matriz de 10 kV, los cut-out, los pararrayos y los aisladores pasatapas primarios del transformador se ejecutará con conductor de cobre recocido temple blando forrado o desnudo de $25\text{ mm}^2$, formando cuellos muertos sin tensiones mecánicas.
\end{itemize}

\section{MONTAJE DEL SISTEMA DE PUESTA A TIERRA}

\begin{enumerate}
  \item \textbf{Excavación del Pozo:} Se excavará un pozo cilíndrico de $0.80\text{ m}$ de diámetro por $2.80\text{ m}$ de profundidad para cada electrodo.
  \item \textbf{Hincado de Varilla:} Se clavará verticalmente en el centro del pozo una varilla de cobre Copperweld de $5/8'' \times 2.40\text{ m}$.
  \item \textbf{Tratamiento con Bentonita Sódica:} Se rellenará el pozo con capas de tierra vegetal cernida mezclada con bentonita sódica (mínimo 2 sacos de 25 kg por pozo) disuelta en abundante agua, compactando firmemente con pisón.
  \item \textbf{Conexión:} La unión entre el electrodo y el conductor bajante de cobre de $35\text{ mm}^2$ se efectuará con conector de bronce tipo AB o soldadura exotérmica (cadweld).
  \item \textbf{Caja de Registro:} Se colocará caja prefabricada en la boca del pozo a nivel de rasante para permitir la medición periódica con telurómetro.
\end{enumerate}

\section{MONTAJE DE LUMINARIAS DE ALUMBRADO PÚBLICO}

\begin{itemize}
  \item Los pastorales de tubo de acero galvanizado de $1 1/2''$ se sujetarán al poste de concreto con dos abrazaderas partidas de dos pernos a una altura de montaje libre de $7.0\text{ m}$ sobre calzada.
  \item La luminaria LED de 50 W se fijará a la punta del pastoral asegurando un apriete uniforme de los pernos de fijación y verificando la orientación del haz lumínico hacia el centro de la calzada.
  \item La conexión eléctrica al conductor de AP ($1 \times 16\text{ mm}^2$) y neutro se realizará mediante conectores de perforación de aislamiento (conector bimetálico tipo diente) estancos con gel sellante impermeable.
\end{itemize}

\section{PRUEBAS, PROTOCOLOS DE CONTROL Y PUESTA EN SERVICIO}

Previo a la energización oficial y entrega a la entidad operadora, se ejecutarán obligatoriamente las siguientes pruebas protocolizadas en presencia de la Supervisión de Obra:

\begin{table}[H]
\centering
\small
\begin{tabular}{|l|l|l|c|}
\hline
\textbf{Prueba Electromecánica} & \textbf{Instrumento Utilizado} & \textbf{Criterio de Aceptación} & \textbf{Norma} \\
\hline
Resistencia de aislamiento MT & Megóhmetro $2,500\text{ V DC}$ & $R_{\text{aisl}} \ge 1,000\ \text{M}\Omega$ a $20^\circ\text{C}$ & IEC 60076 \\
Resistencia de aislamiento BT & Megóhmetro $1,000\text{ V DC}$ & $R_{\text{aisl}} \ge 10.0\ \text{M}\Omega$ & CNE 2011 \\
Resistencia de pozos a tierra & Telurómetro digital 3 bornes & $R_{\text{SED}} \le 10.0\ \Omega$ / $R_{\text{BT}} \le 25.0\ \Omega$ & CNE 2011 \\
Continuidad y secuencia de fases & Fasímetro digital & Rotación de fases directa (R-S-T) & IEC 60034 \\
Rigidez dieléctrica del aceite & Chispómetro homologado & Tensión de ruptura $\ge 35\text{ kV} / 2.5\text{ mm}$ & ASTM D877 \\
Tensión en vacío del transformador & Voltímetro digital RMS & $380 / 220\text{ V} \pm 2.5\%$ en taps nominales & IEC 60076 \\
Prueba de encendido y fotocélula & Tapado manual de sensor & Conmutación correcta de contactores & DGE AP \\
\hline
\end{tabular}
\caption{Protocolos de pruebas oficiales previas a la puesta en servicio.}
\end{table}

Una vez completadas y suscritas las actas de pruebas con resultados conformes, se procederá a la energización en vacío de la línea de 10 kV, seguida de los transformadores y tableros de baja tensión, manteniendo la red en observación durante 48 horas continuas antes de recibir la carga comercial y domiciliaria definitiva.
